\chapter{\MakeUppercase{System Design}}

\section{System Architecture}
The simulator is organized into independent modules. The environment controls time, node creation, message generation, mobility updates, contact discovery, delivery checking, expiration, and metric computation. The routing package contains protocol-specific forwarding logic. This structure allows the same environment to execute different routers without changing the test conditions.

\begin{figure}[h!]
    \centering
    \fbox{\parbox{0.88\textwidth}{\centering
    Image placeholder: add a screenshot or diagram showing the simulator architecture with Environment, Mobility, Node, Message, Routing Protocols, Metrics, CSV Output, and Dashboard.}}
    \caption{System architecture of the DTN routing simulator}
    \label{fig:architecture_placeholder}
\end{figure}

\section{Detailed Design}
The main execution loop performs the following sequence for each time step:

\begin{enumerate}
	\item Generate new messages according to the traffic level.
	\item Update the mobility state of every node.
	\item Detect all node pairs within transmission range.
	\item Invoke the selected router for each contact.
	\item Check whether any message has reached its destination.
	\item Expire messages whose \ac{ttl} has been exceeded.
\end{enumerate}

This design keeps routing decisions local to contacts. A router does not receive global future knowledge; it can only use the state maintained by the nodes and the information available during encounters.

\section{Module Responsibilities and Data Flow}
Each module in the simulator has a narrow responsibility so that routing behavior can be studied without unnecessary coupling. The environment is responsible for global progression of simulation time and for enforcing common rules such as message expiration, delivery accounting, and result aggregation. Nodes are responsible for maintaining local state, including mobility information, buffers, encounter history, and zone memory. Messages carry only the fields needed for routing and evaluation, such as source, destination, age, criticality, and remaining copies.

The data flow between modules follows the same sequence in every run. The environment creates nodes and initializes the selected router. Messages generated during a time step are inserted into source buffers. Mobility updates change node positions, which in turn determine the contact pairs for that step. For each contact, the router reads the local state of both nodes, updates encounter-related information, and applies forwarding rules. After routing exchanges, the environment checks whether any copies have reached their destinations and updates the metric counters. This repeated structure makes protocol comparisons transparent because the only varying component is the router logic.

\section{Proposed Spatio-Semantic Routing Design}
The Spatio-Semantic router computes a utility score for each possible carrier-destination pair. The score combines spatial proximity, direct encounter history, zone co-location, semantic role, and transitive encounter evidence. The composite utility is:

\begin{equation}
	U = 0.18P + 0.28E + 0.20Z + 0.14R + 0.20T
\end{equation}

where \ac{proximity} is normalized spatial proximity, \ac{encounter} is normalized direct encounter frequency, \ac{zone} is normalized visit frequency to the destination zone, \ac{role} is the semantic role bonus, and \ac{transitive} is the two-hop encounter-chain score.

\subsection{Spatial Awareness}
The simulator area is divided into a 5 by 5 grid of zones. Each node records how frequently it visits each zone, with gradual decay over time. A carrier that often visits the destination's zone is considered more useful than a carrier that only happens to be close at the current instant. This is important in \acp{dtn}, where future movement can be more valuable than present distance.

\subsection{Semantic Awareness}
The protocol assigns role bonuses to operationally useful nodes. Drones receive the highest role bonus because they move quickly and can bridge separated regions. Responders receive a smaller bonus because they have higher mobility and operational relevance. Shelters receive a limited bonus because they are meaningful fixed points, especially for messages headed nearby. Civilians receive no additional bonus.

\subsection{Encounter and Transitive Awareness}
Direct encounter history records how often two nodes meet. This captures repeated contacts and stable mobility patterns. Transitive awareness looks at the carrier's strongest peers and checks whether those peers have encountered the destination. This concept is related to history-and-transitivity routing in \ac{prophet}, but here it is one component of a broader spatial and semantic utility function \citep{lindgren2012prophet}.

\subsection{State Aging and Adaptation}
Both encounter history and zone history use decay so that old observations do not dominate recent behavior. This is important in a disaster scenario because mobility patterns can change as people move toward shelters, responders change assignments, or drones are redirected. A node that was useful earlier may become less useful later. Decay therefore gives the protocol limited memory rather than permanent belief.

This adaptive state design also supports fairness in relay selection. Without aging, nodes that were initially central could continue to receive forwarding preference even after their usefulness declines. By gradually reducing stale evidence, the router remains responsive to current movement and contact opportunities. The result is a more realistic approximation of how opportunistic routing should behave in a changing environment.

\section{Forwarding Policy}
The forwarding policy changes according to message priority and copy count. Critical messages are processed before non-critical messages. When many copies remain, the router behaves more aggressively and distributes multiple copies. When fewer copies remain, the router becomes more selective and forwards only when the receiver has better utility. When only one copy remains, handoff is allowed only to a clearly better carrier, with a more permissive rule for drones carrying critical messages.

\begin{algorithm}[h!]
\caption{Spatio-Semantic forwarding decision}
\begin{algorithmic}[1]
\ForAll{messages in sender buffer, ordered as critical first}
	\If{receiver already has the message or no copies remain}
		\State skip the message
	\EndIf
	\If{receiver is the destination}
		\State deliver one copy directly
	\ElsIf{message is critical and copies are above spray threshold}
		\State give a larger share of copies, especially to drones
	\ElsIf{copies are greater than one}
		\State compare receiver utility with sender utility
		\If{utility gain is sufficient}
			\State transfer a proportional number of copies
		\EndIf
	\ElsIf{one copy remains}
		\State hand off only to a clearly better carrier
	\EndIf
\EndFor
\end{algorithmic}
\end{algorithm}

\section{Buffer Management}
The proposed router reserves 35 percent of each buffer for critical messages. Non-critical messages can use only the remaining buffer capacity. If a critical message arrives when the node is full, the router attempts to evict a non-critical message, prioritizing messages with redundant remaining copies. This design directly supports the thesis goal: emergency messages should survive congestion better than routine messages.

\section{Rationale for the Hybrid Design}
The proposed design is hybrid because no single forwarding behavior is sufficient across all traffic conditions. When many copies remain, cautious forwarding wastes opportunity. When only one or two copies remain, aggressive spraying wastes resources and may move the final copy to a poorer carrier. The protocol therefore changes its behavior across the life cycle of a message. This is a deliberate design decision rather than a heuristic convenience.

The same logic applies to the combination of spatial and semantic features. Proximity alone can be misleading in \acp{dtn} because a nearby node may soon move away from the destination area. Role bonuses alone would be too coarse because not every drone or responder is equally useful at every moment. Encounter history alone would underuse static but important destinations such as shelters. The hybrid utility function addresses this by allowing different forms of partial evidence to reinforce one another.

From a system-design perspective, the protocol aims to improve decision quality without making the router overly complex. The maintained state is compact, the scoring function is interpretable, and the forwarding thresholds can be tuned explicitly. This makes the design suitable for educational simulation and provides a clear basis for future extensions such as adaptive weighting or richer semantic roles.

\section{Design Alternatives Considered}
An alternative design would have been to use only prediction-based forwarding with encounter history and transitivity. That approach would be simpler, but it would ignore the physical structure of the environment and the operational meaning of different carriers. Another alternative would have been to use a purely geographic design. That would capture physical position but would miss the fact that some roles, especially drones and responders, are useful not only because of where they are now but because of how they move and what they represent in the scenario.

The final design selected in this thesis balances these alternatives. It preserves the simplicity of a contact-based routing engine while adding enough contextual information to improve the quality of forwarding decisions. This balance is important because a design that is too simple may miss the disaster-specific opportunity, while a design that is too elaborate may become difficult to analyze, justify, and implement within a thesis project.
